\section{An independent search for real defects}
\label{sec:defects}

Since the argument of \N\ cannot establish invalidity, one may still ask whether the release is invalid for some other reason.
Proposition~\ref{prop:exits} lists the possibilities, and each can be checked mechanically. This section does not depend on \N.

\subsection{Exit (i): axioms, \texttt{sorry}, trusting the compiler}

A single-pass scan of all 175,187 source files gives: 0 \lean{axiom} declarations; 0 occurrences of \lean{native_decide}, \lean{ofReduceBool}, \lean{implemented_by}, \lean{extern}, \lean{unsafe}, \lean{opaque}, \lean{csimp};
5 occurrences of \lean{sorry}, all in \path{src/Challenge.lean}, a file added by the community for the comparison tool, absent from the author's archive and not imported by \path{Main.lean}.
The kernel export gives the same answer without relying on text: the only axiom records are \lean{propext}, \lean{Quot.sound} and \lean{Classical.choice}; no record is flagged unsafe or partial;
the names \lean{sorryAx}, \lean{Lean.ofReduceBool}, \lean{Lean.ofReduceNat} and \lean{Lean.trustCompiler} do not occur; all 91 opaque constants are library declarations.
This agrees with \N\ (\S12) and with the first-round hygiene scan of agent~D.

\subsection{Exit (ii): a substituted statement}
\label{sec:rebuild}

The five theorems in \path{Main.lean} have no hypotheses. No declaration of the release is named \lean{RiemannHypothesis}, \lean{riemannZeta} or \lean{DirichletCharacter}.
All 1,695 \lean{notation} commands in the release are \lean{local} and one is \lean{scoped}; there is no global notation and no global macro other than tactic macros, so the parsing of \path{Main.lean} cannot have been altered.
Mathlib is pinned at the official tag \texttt{v4.33.1} with no path override.

The strongest check is a reconstruction. The comparison of a proof export with a statement export is only as good as the statement export, so we rebuilt it:
official Lean 4.33.1 binaries; a project containing only the five statements with \lean{sorry} and \texttt{import Mathlib}; 8,690 precompiled files from the official Mathlib cache; \texttt{lean4export} at tag \texttt{v4.33.0}.

\begin{verdictbox}
The rebuilt file has 350,862,311 bytes and SHA-256 \texttt{bd484ac9…bdd6ab2b9}. It is \textbf{byte-identical} to the published \path{challenge.ndjson}.
The reference against which the proof is compared is therefore determined by official Mathlib and the five statements alone.
\end{verdictbox}

\subsection{Exit (iii): the kernel}
\label{sec:conleche}

The release contains metaprogramming that deserves a look, and this is where a real exploit would have to hide.
The files \path{BuildAid/D1..D5.lean} call \lean{Lean.modifyEnv}, but only to extend a deny-list for library suggestions.
A command \lean{cpt_copy_thms} copies declarations in bulk through the standard \lean{Lean.addDecl}, so each copy passes the kernel.
Two custom tactics, \lean{kexactB} (234 uses) and \lean{kcopy}, assign a term to the goal without elaborator-side type checking; \lean{kcopy} builds raw projection terms:

\excerpt{kcopy}

This is legitimate engineering, but it means that those proofs are guarded by the kernel alone. If a kernel had a defect in the checking of projections, this is where it could be exploited.
Reading the source cannot settle that; an independent replay can.

We built \texttt{con-leche}~\cite{conleche}, a checker with its own term representation and a formalized consistency argument, from source at commit \texttt{bc667610} and ran
\begin{Verbatim}[fontsize=\small,frame=single,rulecolor=\color{gray!50}]
con-leche --verified --compare challenge.ndjson solution.ndjson
\end{Verbatim}
\noindent on the complete 61.57\,GB export. The run ended with exit code 0 and this output:
\VerbatimInput[fontsize=\scriptsize,frame=single,rulecolor=\color{gray!50},breaklines=true,breakanywhere=true]{../evidence/con-leche/compare.stdout.log}
\noindent All 1,468,418 declarations of the proof export were accepted, and each of the five statements of the challenge export was found to be justified by a constant of the same name and type whose statement reads the same declarations on both sides.
The run took 73 minutes on 16 cores with a peak of 47.4\,GB of memory. The verification page~\cite{mit} reports the same counts for \texttt{con-leche} and passes by further independent kernels, which we did not re-run.

\subsection{Exit (iv): source versus checked artefact}

The hashes in Table~\ref{tab:prov} tie the archive, the mirror and the two exports to the published values.
The one link we did not reproduce is the build that turns the 175,187 source files into the proof export; for it we rely on the sandboxed rebuild reported on the verification page.
This limitation affects the claim that the \emph{source} proves the theorem. It does not affect the claim that the \emph{export} does: the export was checked directly, and the dependency analysis of this note was performed on it.

\subsection{Result}

None of the four exits shows any sign of being open.
